\chapter{Methodology}

\begin{chapterintro}
This chapter details the research method of enquiry: Participatory Co-design. We discuss the participation and activities in each research cycle.
\end{chapterintro}

\section{Tell-Make-Enact Co-design Cycle}

\begin{figure}[h]
\centering
\begin{tikzpicture}[
    oval/.style={
        draw=lightgray!60!black,
        line width=2.5pt,
        ellipse,
        minimum width=7cm,
        minimum height=3cm,
        align=center,
        font=\bfseries\LARGE
    },
    arrow/.style={
        <->,
        >=Latex,
        ultra thick,
        black
    }
]

% Nodes (positioned like your image)
\node[oval] (make) at (0,3) {Make};
\node[oval] (tell) at (-5,0) {Tell};
\node[oval] (enact) at (5,0) {Enact};

% Curved arrows (cycle flow)
\draw[arrow] (tell) to[bend left=30] (make);
\draw[arrow] (make) to[bend left=30] (enact);
\draw[arrow] (enact) to[bend left=30] (tell);

\end{tikzpicture}
\caption{The Tell--Make--Enact Co-design Cycle}
\label{fig:tell-make-enact}
\end{figure}



The research method adopted for this research is Participatory Design (PD) as has been discussed under the sub-heading: “Participatory Design for Development” making use of the tell-make-enact cycle over three iterations. The cycle was coined by Brandt and others \cite{simonsen2013}. Brandt suggests that the traditional participatory action research design cycle of Plan $\rightarrow$ Do $\rightleftarrows$ Observe $\rightarrow$ Reflect is replaced by a Tell--Make--Enact cycle,
that allows prototyping to happen in any stage and acknowledges that design entails movement in both directions around the cycle. Thus, we use the PD methodology with a change in the design cycle to allow for fluid Co-design (designing together on par) with potential users.


\bigskip

\textbf{Tell $\rightarrow$ Make $\rightarrow$ Enact}

\begin{itemize}
    \item \textbf{Tell:} A Co-designer describes an employment-related idea
    \item \textbf{Make:} A prototype of the idea is created
    \item \textbf{Enact:} The prototype is used to generate new ideas and insights
\end{itemize}

\bigskip

Our research focused on trainers and students who were currently taking or had recently graduated from job-readiness programs. In part, this is because we believe that it is important to work with NGOs as gatekeepers in ICT4D projects since they already have programs and relationships with the target population. As a result, they can help provide additional insight. Intermediaries are often trusted and sometimes able to bridge language gaps. In addition, it is difficult to make ICT interventions work in absence of any support infrastructure. We expect that without JR training and access to job-seeking services through the NGO, the participants would not be able to benefit from a specialized application. Working with two organisations allows us to get a broad understanding of the problem, from multiple different perspectives and obtain a collective solution which is potentially very useful for similar contexts around the world.

The research is conducted using participant observation, trainer interviews, student interviews, student workshops and one-on-one walkthrough sessions. We conducted the participant observation as a long-term source of source of unspoken information as well as tacit knowledge. We conducted interviews and one-on-one walkthrough sessions with both student and trainers. During the interviews and walkthrough sessions, we recorded our participants’ dialogues and later transcribed the recordings. We used both NVivo and manual comparison of findings. For the manual comparisons, data were tabulated into columns in Microsoft Excel, where each column represents the same data attribute. I.e. Separate columns for answers to different questions and then comparing the individuals’ responses and suggestions. Some excel sheets of data analysis have been included in the appendices of this dissertation.

\section{Participation}

In this research, we worked with unemployed persons from two NGOs (not-for-profit-organisations). We chose to work with two separate NGOs for our findings to be rich and greatly generalizable. We worked with proactive unemployed persons, those attending the NGO-based JR programs. The job seekers who are mostly low-skilled are sometimes referred to as day-workers as they perform short-term jobs \cite{chepken2012tele, chepken2011}. Limiting our engagement to willing participants from the JR programs. As such, we eliminated those who were not actively looking for jobs, i.e. those who were not present because of family responsibility and those who were not leveraging the opportunities put out by the JR programs.

We expect that the individuals who have attended the JR programs, which included computer learning would be better positioned to leverage a new ICT-based intervention. The exposure to technology, access to computer and internet facilities and the presences of trainers at the NGO sites would assist them in assimilating new technology into their job search. Thus, any technology developed is being used in the context of existing socio-structural resources.

We worked with the NGOs for a period of 18 months. By spending a long period of time with the NGOs, we got a chance to interact with many unemployed students from different intake groups, thus, broadening our understanding of unemployment and technology. In each cycle of our research, we worked with a new set of participants, particularly because the JR programs were shorter than our research cycles, which meant that our participants would have graduated to enter the job searching world before each subsequent cycle. By having different participants in each cycle, we got design suggestions and other feedback from a larger set of participants.

Table ~\ref{tab:research-activities-cycle} gives details on the activities we conducted in each cycle. In the first, we focused on understanding the problem and in subsequent cycles we sought out ways in which one of the problems could be tacked. Outside the formal research activities, the researcher volunteered at the NGO cites to assist participants with computer training, interview preparation and CV creation and proofreading. The interactions allowed the researcher to understand the bottlenecks in the job search process.


\begin{table}[htbp]
\centering
\caption{Research Activities Listed by Cycle}
\label{research_activities_by_cycles}
\renewcommand{\arraystretch}{1.35}
\begin{tabularx}{\textwidth}{
|>{\centering\arraybackslash}p{2.6cm}|
>{\centering\arraybackslash}p{1.2cm}|
>{\centering\arraybackslash}X|
>{\centering\arraybackslash}p{2.6cm}|
>{\centering\arraybackslash}p{2.6cm}|
}
\hline
\rowcolor{gray!30}
\textbf{Cycle} & \textbf{NGO} & \textbf{Activity} & \textbf{Student Participants} & \textbf{Trainer Participants} \\
\hline

\multirow{4}{=}{\centering\textbf{Cycle 1: Needs Assessment}}
& AT & Workshop & 5 & -- \\
\cline{2-5}
& AT & Interviews & 2 & 2 \\
\cline{2-5}
& TZN & Workshop & 33 & -- \\
\cline{2-5}
& TZN & Interviews & 2 & 2 \\
\hline

\multirow{2}{=}{\centering\textbf{Cycle 2: Ideas Focus}}
& AT & Walkthroughs & 6 & 2 \\
\cline{2-5}
& TZN & Walkthroughs & 5 & 1 \\
\hline

\multirow{2}{=}{\centering\textbf{Cycle 3: Deployment}}
& AT & Deployment and Walkthrough & 5 & -- \\
\cline{2-5}
& TZN & Deployment and Walkthrough & 5 (1 incomplete) & -- \\
\hline

\end{tabularx}
\label{tab:research-activities-cycle}
\end{table}




\section{Ethics}

\subsection{Testing a Potential Solution}

Before embarking on any research, the researcher should analyse whether there is a problem or there is an opportunity for improvement. Before the beginning of this project in 2016, the researcher found out that the Unemployment Rate in South Africa was over 25\% in 2016 \cite{statssa2016} and thus, the opportunity to assist in bringing down to the government’s goal of 6\% \cite{statssa2015, npc2013}. 

To understand the problem, we decided to partner with NGOs who focused on training job-seekers who were struggling to get jobs. Through the NGOs, we were able to engage over 50 participants in our project at various stages of the project and to meet over many other job seekers during the sessions when the researcher volunteered to offer the job-seekers, Computer skills (typing, online search, emailing, etc.). 

It was a valid question to ask, should a job-seeker with no source of income be engaged in a project which’s value is in the future? However, we considered it plausible to continue the research as we expected the future value to negate the time taken to take part in the research. However, at the beginning of each session with the job seekers, we let them know that they were free to stop participating at any point and during one of our sessions, one of our workshop participants excused himself to go to his part-time job.

%================================================
\subsection{Hidden Identity and Hidden Aptitude}

Racial segregation, discrimination major unethical issues that occur frequently in society, however, with ICTs individuals are able to avert being segregated because of their skin colour or race as ICTs are capable of hiding identities \cite{chepken2012tele}. 

The usage of hidden identities is positive to avoid discrimination however, ICTs go to the extent of hiding people’s aptitudes as computers and other smart devices perform work that is supposed to be performed by individuals which bring an issue around whether it is fair to hide individuals’ abilities as opposed to expose the potential of the individual (who is in this instance a technology user). 

In job-seeking for example: should an application hide the potential of a job seeker to an employer seeking the best candidate for a task? The issue of hidden aptitude is particularly relevant with blue-collar jobs where an employer has a limited time to test an individual’s abilities as jobs last short periods: hours, weeks and at most, a year where probations do not apply \cite{chepken2012tele}. 

The proposed ‘Employment Service Bill’ of South Africa aims to promote the employment of disabled people \cite{gitau2012}. However, the same ICTs which would hide aptitude would work differently for the disabled and other home-bound people, because it would eliminate tasks that they are no longer able to perform by making use of a different/non-impaired sense \cite{stefano2016} or allow working from a remote area \cite{stefano2016}. However, when individuals work at different geographical locations, it may mean that some will have to work at night or during unsociable hours \cite{stefano2016}.  
%==================================================
\subsection{Compensation}

Compensation or remuneration given to research subjects for participating in research projects is useful as a form of motivation \cite{zimmerman2011, ssozi2017} in addition to the intrinsic motivation to solve the problem at hand. Similarly to other ICT research, we compensated participants with food \cite{winschiers2009} and airtime \cite{gupta2012}.


In some instances, compensation to the research subjects serves both the researcher and the subjects, for example in rural Uganda participants were given phones for use for the research’s goal as well as personal use and the participants felt they gained social status in the community by having phones which were better than other community members’ phones. In India including users in videos was used as both incentive to motivate their participation and as an aid to the adoption of the tool for others \cite{gandhi2007}, thus the form of compensation helped the effectiveness and sustainability of the project.  

Although we did not experience the bidirectional benefit of compensation, in this research we found that the volunteer work that the research did with the job seekers made some of them more willing and more comfortable to participate in our research, such that most of them volunteered to be part of the research. In cycle 2 when participants were not compensated, the \replaced{researcher}{research} struggled to get enough participants.

\section{Cycle 1: Understanding ICT Use and Unemployment}

In Cycle 1 we focused on understanding unemployment and our unemployed audience. We interviewed two trainers and two student job seekers to get initial insights into unemployment in Cape Town. After the interview, we performed a workshop with more students at each of our partner sites. Table~\ref{tab:research-activities-cycle} gives the details of the participation. The researcher also attended some of the student’s JR classes and assisted them with:

\begin{itemize}
\item editing their CVs
\item learning to use computers
\end{itemize}


%------------------------------------------------

\section{Cycle 2: Ideas Focus}

In Cycle 2 we focused on creating prototypes of the selected intervention and developing the artefact. We performed two sets of walkthroughs, the first was to get inputs on the hi-fidelity prototypes, and the second was to collect input on the usability, utility and feasibility of the working artefact. The improvements, suggestions and other participant inputs from the first set of walkthroughs were used to improve the design of the artefact, which was named Khanya Job Network.

We developed the artefact using PHP and HTML web pages, a MySQL database and an Apache server. All of which were free and accessible. The artefact was developed in the XAMPP framework.

We implemented a system that uses the personal information provided by job seekers to:

\begin{itemize}
\item Offer job seekers tips in finding jobs
\item Generate CVs
\item Generate cover letters
\item Facilitate the generation of customized job opportunities which were sent by email and SMS to job seekers
\end{itemize}


%------------------------------------------------

\section{Cycle 3 Deployment}

In the third cycle, we deployed the artefact for a period of two weeks. At the beginning of the two weeks, we had ten participants to register on Khanya Job Network. During the deployment, our participants received job notifications by SMS, Email or both depending on their choice during registration. At the end of the two weeks, we performed a final walkthrough session to evaluate the application.